% !TeX root = part6.tex
\documentclass[../final.tex]{subfiles}
\begin{document}
\ifSubfilesClassLoaded{\setcounter{chapter}{3}}{}

% Lecture 6, 05.10.2026. Sources: Desktop/11.jpg, 22.jpg, 33.jpg, 44.jpg.
\chapter{Термодинамика бозе-конденсата}
\rsubtitle{Теплоёмкость, давление и изотермы}
\section{Теплоёмкость вблизи конденсации}
\subsection{Число частиц и внутренняя энергия}
Рассматриваем трёхмерный идеальный бозе-газ с сохраняющимся числом частиц.
Температура измеряется в энергетических единицах: $T=k_BT_{\mathrm K}$.
Обозначим $T_c=T_{\text{БЭ}}$, $x=\mu/T$ и
\begin{equation*}
 A=\frac{g(2m)^{3/2}}{4\pi^2\hbar^3},\qquad
 F_\alpha(x)=\sum_{n=1}^{\infty}\frac{e^{nx}}{n^\alpha},\qquad x<0.
\end{equation*}
Здесь $F_\alpha'(x)=F_{\alpha-1}(x)$,
$F_\alpha(0)=\zeta(\alpha)$ при $\alpha>1$.
В нормальной фазе $T>T_c$ число частиц равно
\begin{align*}
 N&=AVT^{3/2}\Gamma(3/2)F_{3/2}(x),\\
 F_{3/2}(x)&=\zeta(3/2)-2\sqrt\pi\sqrt{-x}+O(|x|),\qquad x\to0^-.
\end{align*}
На границе перехода
\begin{equation*}
 N=AVT_c^{3/2}\Gamma(3/2)\zeta(3/2),\qquad
 T_c=\left[\frac{N}{AV\Gamma(3/2)\zeta(3/2)}\right]^{2/3}.
\end{equation*}
Энергия имеет две ветви:
\begin{equation*}
 E=AVT^{5/2}\Gamma(5/2)
 \begin{cases}
  \zeta(5/2),&T\leq T_c,\\
  F_{5/2}(x),&T>T_c.
 \end{cases}
\end{equation*}
При постоянных $N,V$ теплоёмкость равна
\begin{align*}
 C_V^-&=\frac52AVT^{3/2}\Gamma(5/2)\zeta(5/2),\\
 C_V^+&=AV\Gamma(5/2)
 \left[\frac52T^{3/2}F_{5/2}(x)+T^{5/2}F_{3/2}(x)x'\right].
\end{align*}
Штрих означает производную по $T$ при постоянных $N,V$.
Поскольку $x'(T_c^+)=0$, теплоёмкость на границе непрерывна:
$C_V(T_c^+)=C_V(T_c^-)$.

\clearpage
\subsection{Вторая производная отношения $\mu/T$}
Корневая поправка к числу частиц даёт
\begin{align*}
 \frac{N}{AV\Gamma(3/2)T^{3/2}}
 &=\zeta(3/2)-2\sqrt\pi\sqrt{-x}+O(|x|),\\
 -\frac{3N}{2AV\Gamma(3/2)T^{5/2}}
 &\sim\frac{\sqrt\pi}{\sqrt{-x}}x',\\
 x'&\sim-\frac{3N}{2AV\Gamma(3/2)\sqrt\pi\,T^{5/2}}\sqrt{-x}
 \longrightarrow0.
\end{align*}
Положим $B(T)=3N/[2AV\Gamma(3/2)\sqrt\pi\,T^{5/2}]$.
Дифференцируя главный член $x'\sim-B(T)\sqrt{-x}$, находим
\begin{align*}
 x''&\sim-B'(T)\sqrt{-x}+\frac{B(T)}{2\sqrt{-x}}x',\\
 x''(T_c^+)&=-\frac12B^2(T_c)
 =-\frac12\left[\frac{3N}{2AV\Gamma(3/2)\sqrt\pi\,T_c^{5/2}}\right]^2\\
 &=-\frac{9\zeta^2(3/2)}{8\pi T_c^2}.
\end{align*}
Первая добавка исчезает, а во второй отношение $x'/\sqrt{-x}$
стремится к $-B(T_c)$. Поэтому вторая производная имеет конечный предел.

\subsection{Скачок производной теплоёмкости}
В пределе $T\to T_c^+$ члены, содержащие $x'$ и
$F_{1/2}(x)(x')^2$, исчезают. Остаётся
\begin{align*}
 \Delta\!\left(\frac{dC_V}{dT}\right)
 &\equiv C_V'(T_c^+)-C_V'(T_c^-)\\
 &=AVT_c^{5/2}\Gamma(5/2)\zeta(3/2)\,x''(T_c^+)\\
 &=-\frac{27}{16\pi}\frac{N}{T_c}\zeta^2(3/2)
 \simeq-3.666\frac{N}{T_c}.
\end{align*}
Здесь использованы $\Gamma(5/2)=\tfrac32\Gamma(3/2)$ и уравнение для $T_c$.
При этом $C_V(T_c)/N=\tfrac{15}{4}\zeta(5/2)/\zeta(3/2)\simeq1.926$.

\begin{rbox}[unbreakable,left=14pt,right=14pt,top=10pt,bottom=8pt]{[]}
\centering
\captionsetup{font=footnotesize,hypcap=false,skip=5pt}
\begin{tikzpicture}[>=Stealth,font=\small,x=2.4cm,y=1.05cm]
 \draw[->] (0,0)--(4.3,0) node[right] {$T/T_c$};
 \draw[->] (0,0)--(0,2.35) node[above] {$C_V/N$};
 \draw[dashed,red3] (0,1.5)--(4,1.5);
 \node[left] at (0,1.5) {$3/2$};
 \draw[dotted] (1,0)--(1,1.926);\node[below] at (1,0) {$1$};
 \draw[p3,thick,domain=0:1,samples=50] plot (\x,{1.925671675*\x^1.5});
 \draw[p3,thick] plot file {\subfix{../image/bose_cv.dat}};
 \node[above,p3] at (1,1.926) {$1.926$};
\end{tikzpicture}
\captionof{figure}{Теплоёмкость непрерывна, её производная испытывает скачок.
 При высоких температурах $C_V/N\to3/2$.}
\end{rbox}

\clearpage
\section{Давление идеального бозе-газа}
\subsection{Большой термодинамический потенциал}
Для нерелятивистского спектра $\varepsilon=p^2/(2m)$
плотность состояний равна $D(\varepsilon)=AV\sqrt\varepsilon$.
В термодинамическом пределе
\begin{equation*}
 \Omega=AVT\int_0^\infty\sqrt\varepsilon\,
 \ln\!\left(1-e^{(\mu-\varepsilon)/T}\right)d\varepsilon.
\end{equation*}
Интегрирование по частям с учётом
$d\ln(1-e^{(\mu-\varepsilon)/T})/d\varepsilon
 =[T(e^{(\varepsilon-\mu)/T}-1)]^{-1}$ даёт
\begin{equation*}
 \Omega=-\frac23AV\int_0^\infty
 \frac{\varepsilon^{3/2}\,d\varepsilon}{e^{(\varepsilon-\mu)/T}-1}
 =-\frac23E=-pV.
\end{equation*}
Следовательно,
\begin{equation*}
 p=\frac{2E}{3V}=\frac23AT^{5/2}\Gamma(5/2)F_{5/2}(\mu/T),\qquad T>T_c.
\end{equation*}

\subsection{Давление при наличии конденсата}
При $T\leq T_c$ химический потенциал равен нулю и
\begin{align*}
 E&=AVT^{5/2}\Gamma(5/2)\zeta(5/2),\\
 p_{\mathrm{sat}}(T)&=\frac23AT^{5/2}\Gamma(5/2)\zeta(5/2),\\
 N_{\mathrm{ex}}&=AVT^{3/2}\Gamma(3/2)\zeta(3/2),\qquad
 N_0=N-N_{\mathrm{ex}}.
\end{align*}
Давление определяется возбуждёнными частицами и при заданной температуре
не зависит от полного числа частиц. Частицы в основном состоянии
с $\varepsilon_0=0$ не вносят вклада в энергию и давление идеального газа.

При фиксированных $N,T$ критический объём равен
\begin{equation*}
 V_c(T)=\frac{N}{AT^{3/2}\Gamma(3/2)\zeta(3/2)}.
\end{equation*}
При сжатии ниже $V_c$ число возбуждённых частиц уменьшается пропорционально $V$,
доля конденсата растёт, а давление остаётся равным $p_{\mathrm{sat}}(T)$.

\clearpage
\section{Изотермы квантовых газов}
\subsection{Давление и химический потенциал}
При $V>V_c$ бозе-газ находится в нормальной фазе: $\mu<0$.
Уравнение для числа частиц определяет химический потенциал:
\begin{equation*}
 F_{3/2}(\mu/T)=\zeta(3/2)\frac{V_c}{V}.
\end{equation*}
При $V\leq V_c$ имеется конденсат и $\mu=0$.
Сжатие в этой области переводит частицы в основное состояние.

\begin{rbox}[unbreakable,left=16pt,right=16pt,top=12pt,bottom=9pt]{[]}
\centering
\captionsetup{font=footnotesize,hypcap=false,skip=6pt}
\begin{tikzpicture}[>=Stealth,font=\small,x=2.15cm,y=0.95cm]
 \draw[->] (0,0)--(5.15,0) node[right] {$V/V_c$};
 \draw[->] (0,0)--(0,4.65) node[above] {$p/p_{\mathrm{sat}}$};
 \draw[dotted] (1,0)--(1,4);\node[below] at (1,0) {$1$};
 \node[left] at (0,1) {$1$};
 \begin{scope}
  \clip (0,0) rectangle (5,4.45);
  \draw[red3,thick] plot file {\subfix{../image/lecture6_fermi_isotherm.dat}};
  \draw[dashed,domain=0.45:5,samples=90] plot (\x,{1.94737/\x});
  \draw[p3,very thick] (0,1)--(1,1);
  \draw[p3,thick] plot file {\subfix{../image/lecture6_bose_isotherm.dat}};
 \end{scope}
 \node[red3,right] at (1.65,2.95) {ферми-газ};
 \node[right] at (2.65,1.4) {классический газ};
 \node[p3,right] at (2.65,0.5) {бозе-газ};
 \node[p3,above] at (0.52,1.08) {конденсат};
 \begin{scope}[yshift=-3.25cm,y=1.55cm]
  \draw[->] (0,0)--(5.15,0) node[right] {$V/V_c$};
  \draw[->] (0,-1.2)--(0,0.4) node[above] {$\mu/T$};
  \draw[p3,very thick] (0.05,0)--(1,0);
  \draw[p3,thick] plot file {\subfix{../image/lecture6_bose_mu_volume.dat}};
  \draw[dotted] (1,0)--(1,-1.1);\node[above] at (1,0) {$1$};
  \node[p3,above] at (0.5,0.08) {$\mu=0$};
  \node[p3] at (2,-0.9) {$\mu<0$};
 \end{scope}
\end{tikzpicture}
\captionof{figure}{Изотермы при одинаковых $N,T,m,g$ и химический потенциал
 бозе-газа. Объём нормирован на $V_c(T)$, давление --- на $p_{\mathrm{sat}}(T)$.}
\end{rbox}

При больших объёмах все три газа переходят в классический режим $pV=NT$.
При сжатии ферми-газа давление растёт сильнее из-за принципа Паули;
в сильно вырожденной области $p\propto(N/V)^{5/3}$.
У бозе-газа давление ниже классического и выходит на горизонтальный участок.

\clearpage
\section{Сравнение с газом Ван-дер-Ваальса}
Для взаимодействующих частиц фазовый переход газ--жидкость описывается,
в частности, уравнением Ван-дер-Ваальса
\begin{equation*}
 \left(p+a\frac{N^2}{V^2}\right)(V-Nb)=NT.
\end{equation*}
Параметры $a,b>0$ учитывают притяжение и собственный объём частиц.
Ниже критической температуры формальная изотерма содержит участок
с $\partial p/\partial V>0$, соответствующий механической неустойчивости.
В равновесии петля заменяется горизонтальным участком сосуществования
жидкости и пара:
\begin{equation*}
 p(V_\ell,T)=p(V_g,T)=p_{\mathrm{coex}}(T),\qquad
 \int_{V_\ell}^{V_g}[p(V,T)-p_{\mathrm{coex}}(T)]\,dV=0.
\end{equation*}
Последнее равенство --- правило равных площадей Максвелла.

\begin{rbox}[unbreakable,left=16pt,right=16pt,top=12pt,bottom=9pt]{[]}
\centering
\captionsetup{font=footnotesize,hypcap=false,skip=6pt}
\begin{tikzpicture}[>=Stealth,font=\small,x=2.4cm,y=2.4cm]
 \draw[->] (0,0)--(4.85,0) node[right] {$V/V_{\text{кр}}$};
 \draw[->] (0,0)--(0,1.5) node[above] {$p/p_{\text{кр}}$};
 \draw[dashed] plot file {\subfix{../image/lecture6_vdw_isotherm.dat}};
 \draw[p3,very thick] (0.603402,0.646998)--(2.348842,0.646998);
 \draw[dotted] (0.603402,0)--(0.603402,0.646998);
 \draw[dotted] (2.348842,0)--(2.348842,0.646998);
 \node[below] at (0.603402,0) {$V_\ell$};
 \node[below] at (2.348842,0) {$V_g$};
 \node[p3,above] at (1.5,0.67) {жидкость и пар};
 \node[right] at (2.8,1.12) {$T=0.9T_{\text{кр}}$};
 \begin{scope}[xshift=2cm,yshift=-4.7cm,x=2cm,y=1.4cm]
  \draw[->] (0,0)--(3.5,0) node[right] {$T$};
  \draw[->] (0,0)--(0,2.35) node[above] {$p$};
  \draw[p3,thick] (0.45,0.25) .. controls (1.2,0.5) and (2,1.3) .. (2.6,1.95);
  \fill[p3] (2.6,1.95) circle (1.5pt);
  \node[right] at (2.6,1.95) {критическая точка};
  \node at (0.95,1.25) {жидкость};
  \node at (2.4,0.45) {газ};
 \end{scope}
\end{tikzpicture}
\captionof{figure}{Изотерма Ван-дер-Ваальса с участком сосуществования фаз
 и схема линии жидкость--газ на плоскости $(T,p)$.}
\end{rbox}

Горизонтальные участки внешне сходны, но имеют разный смысл.
Для газа Ван-дер-Ваальса сосуществуют пространственные фазы с разными плотностями;
для идеального бозе-газа происходит перераспределение частиц между основным
и возбуждёнными состояниями.

\clearpage
\section{Экспериментальная бозе-конденсация}
\subsection{Разреженные атомные газы}
В 1995 году в JILA, совместной лаборатории Университета Колорадо и NIST,
группа Эрика Корнелла и Карла Вимана получила конденсат атомов $^{87}\mathrm{Rb}$.
Независимо группа Вольфганга Кеттерле в MIT получила конденсат атомов натрия.
За получение конденсатов в разреженных газах щелочных атомов и исследование
их свойств Корнелл, Виман и Кеттерле получили Нобелевскую премию по физике 2001 года.

В опыте с рубидием конденсат появлялся при температуре около
$170\,\text{нК}$ и плотности порядка $2.5\cdot10^{12}\,\text{см}^{-3}$.
При дальнейшем охлаждении был получен почти чистый конденсат примерно
из $2000$ атомов при температуре порядка $20\,\text{нК}$.
Атомы удерживались в ловушке; поэтому пространственное распределение плотности
отличалось от однородной модели газа в ящике.

\subsection{Зависимость критической температуры от массы}
Для однородного идеального газа
\begin{equation*}
 T_{c,\mathrm K}=\frac{2\pi\hbar^2}{mk_B}
 \left[\frac{n}{g\zeta(3/2)}\right]^{2/3},\qquad n=\frac NV.
\end{equation*}
При фиксированных $n,g$ коэффициент $A\propto m^{3/2}$, откуда
\begin{equation*}
 T_c^{3/2}\propto\frac1A\propto m^{-3/2},\qquad T_c\propto\frac1m.
\end{equation*}
Чем меньше масса частицы, тем выше температура квантового вырождения
при одной и той же плотности.

\subsection{Оценка для гелия-4}
Для $^4\mathrm{He}$ имеем $g=1$, $m=M_r m_u$,
где $M_r\simeq4.0026$, а $m_u$ --- атомная единица массы.
Введём молярный объём $v_m=V/\nu$; тогда $n=N_A/v_m$.
Если численное значение $v_m$ задано в $\text{см}^3/\text{моль}$, то
\begin{equation*}
 T_{c,\mathrm K}\simeq
 \frac{114.6\,\mathrm K}
 {M_r\,[v_m/(\text{см}^3/\text{моль})]^{2/3}}.
\end{equation*}
При $v_m\simeq27.6\,\text{см}^3/\text{моль}$ идеальная модель даёт
$T_{c,\mathrm K}\simeq3.1\,\mathrm K$.
Температура сверхтекучего перехода жидкого $^4\mathrm{He}$ при давлении
насыщенного пара составляет $T_\lambda\simeq2.17\,\mathrm K$.
Жидкий гелий является взаимодействующей системой, поэтому формула идеального
газа даёт лишь оценку температурного масштаба.

\end{document}
